\section*{Step 254: Selberg Orthogonality Carrier}

\paragraph{Carrier.}
Let \(L_1,L_2\) be distinct primitive elements of the Selberg class with prime
coefficients \(a_p(L_i)\).  Selberg orthogonality predicts
\[
\frac{1}{\log\log x}
\sum_{p\le x}\frac{a_p(L_1)\overline{a_p(L_2)}}{p}\longrightarrow 0.
\]
The associated residual is
\[
\Xi_{\mathrm{SOC}}=
\limsup_{x\to\infty}
\left|
\frac{1}{\log\log x}
\sum_{p\le x}\frac{a_p(L_1)\overline{a_p(L_2)}}{p}
\right|.
\]

\paragraph{CRE audit.}
SOC is not equivalent to the Riemann Hypothesis.  It is a cross-correlation
statement about coefficient independence between distinct primitive
Selberg-class L-functions.  It is also not identical to an individual GRH
statement for one \(L\)-function.  GRH-type assumptions together with
Selberg-class structural decomposition are a natural conditional route to SOC,
but this is not an audited equivalence.

\paragraph{Classification.}
SOC lies outside the Riemann-RH Dichotomy.  It also lies outside the
single-\(L\)-function Selberg-class Dichotomy generalization from step 232,
because its residual is defined between two primitive \(L\)-functions.

\paragraph{Extension suggested.}
SOC suggests a separate Selberg-class cross-correlation extension: typed
carriers whose residuals measure coefficient decorrelation between primitive
objects rather than the RH analogue of one object.

\paragraph{Verdict.}
\[
\mathrm{V\_selberg\_orthogonality\_outside\_dichotomy}.
\]
No proof of SOC, GRH, or RH is claimed.
